\chapter{Exploratory Data Analysis}
\label{ch:eda}

Until now we have built probability models for random phenomena. When we meet a
new phenomenon, though, we usually start from the other end: we collect a set of
observations—a \emph{dataset}—and try to learn what kind of randomness produced
it. \emph{Exploratory data analysis} (EDA) is the art of summarizing a dataset,
graphically and numerically, so that its essential features stand out. These
summaries are \emph{descriptive}: they let the data speak for themselves,
without yet inferring anything about an underlying distribution.

Throughout this chapter we use as a running example the durations (in seconds) of
$272$ eruptions of the Old Faithful geyser. Simply listing the numbers, or even
sorting them, reveals little; we need better summaries.

\section{Graphical Summaries}

\subsection{The histogram}

A \emph{histogram} divides the range of the data into intervals, or \emph{bins}
$B_1,\dots,B_m$, and draws a bar over each. We scale the histogram so the total
area is one, which makes it a crude estimate of an underlying density: the height
over bin $B_i$ is
\[
  \text{height} = \frac{\text{number of observations in }B_i}{n\,|B_i|},
\]
where $|B_i|$ is the bin width. The area of each bar then equals the proportion of
data in that bin.

The choice of bin width is a genuine tradeoff. Too narrow, and the histogram
dissolves into a jagged row of isolated spikes; too wide, and real structure is
smoothed away. Figure~\ref{fig:histbw} shows the geyser data at three bin widths.
The middle choice reveals the dataset's most important feature—\emph{two}
clusters of eruption durations, near $120$ and $270$ seconds—that the widest bins
erase entirely.

\begin{figure}[ht]
  \centering
  \includegraphics[width=\linewidth]{hist_binwidth}
  \caption{The same geyser data at three bin widths. Too small produces noise;
  too large hides the two modes. In practice one chooses the bin width by trial
  and error until the picture looks reasonable.}
  \label{fig:histbw}
\end{figure}

\subsection{Kernel density estimates}

A histogram is simple but jagged, and sensitive to where the bin edges fall. A
\emph{kernel density estimate} (KDE) produces a smooth curve instead. The idea is
to ``place a small pile of sand'' over each observation and add the piles up;
where data accumulate, the sand piles high. Formally one chooses a \emph{kernel}
$K$—a symmetric density giving each pile its shape—and a \emph{bandwidth} $h$
controlling each pile's width.

\begin{figure}[ht]
  \centering
  \includegraphics[width=0.62\linewidth]{kde_faithful}
  \caption{A kernel density estimate of the geyser data overlaid on a histogram.
  The two modes are now unmistakable and the curve is smooth.}
  \label{fig:kdefaithful}
\end{figure}

A kernel $K$ typically satisfies $K(u)\ge 0$, integrates to one, is symmetric
about zero, and vanishes outside $[-1,1]$. Common choices include the
Epanechnikov kernel $K(u)=\tfrac34(1-u^2)$ and the triweight kernel
$K(u)=\tfrac{35}{32}(1-u^2)^3$; the normal kernel is also popular, though it
violates the last condition. Figure~\ref{fig:kernels} compares several.

\begin{figure}[ht]
  \centering
  \includegraphics[width=0.6\linewidth]{kernels}
  \caption{Several standard kernels. The choice of kernel matters far less than
  the choice of bandwidth.}
  \label{fig:kernels}
\end{figure}

The bandwidth plays exactly the role the bin width played for the histogram: too
small gives a spiky curve, too large oversmooths and can wash out the modes
(Figure~\ref{fig:kdebw}). The choice of \emph{kernel}, by contrast, has only a
minor effect on the result.

\begin{figure}[ht]
  \centering
  \includegraphics[width=\linewidth]{kde_bandwidth}
  \caption{The bandwidth $h$ controls smoothness. Two well-known kernels with the
  same bandwidth produce nearly identical curves; the bandwidth is what matters.}
  \label{fig:kdebw}
\end{figure}

\subsection{The empirical distribution function}

A third graphical summary plots the data cumulatively. The \emph{empirical
distribution function} (or empirical cdf) is
\[
  F_n(x)=\frac{\text{number of observations}\le x}{n},
\]
the proportion of the dataset at or below $x$. Its graph is a staircase: flat
below the minimum, equal to one above the maximum, and rising by a jump of
$1/n$ at each observation.

\begin{example}[A small empirical cdf]
For the dataset $\{4,3,9,1,7\}$, the empirical cdf is $0$ for $x<1$, jumps to
$\tfrac15$ at $x=1$, to $\tfrac25$ at $x=3$, to $\tfrac35$ at $x=4$, and so on
up to $1$ for $x\ge 9$ (Figure~\ref{fig:ecdftoy}). The closed and open circles
emphasize that the function takes its \emph{higher} value at each jump point.
\end{example}

\begin{figure}[ht]
  \centering
  \includegraphics[width=0.52\linewidth]{ecdf_toy}
  \caption{The empirical distribution function of a five-point dataset.}
  \label{fig:ecdftoy}
\end{figure}

For the geyser data (Figure~\ref{fig:ecdffaithful}), the empirical cdf rises
\emph{steeply} where observations crowd together—near $120$ and $270$
seconds—and flattens where they thin out. Steepness of $F_n$ is the cumulative
signature of a mode.

\begin{figure}[ht]
  \centering
  \includegraphics[width=0.52\linewidth]{ecdf_faithful}
  \caption{Empirical distribution function of the geyser data; the two steep
  stretches correspond to the two clusters of durations.}
  \label{fig:ecdffaithful}
\end{figure}

\subsection{The scatterplot}

For a \emph{bivariate} dataset of pairs $(x_1,y_1),\dots,(x_n,y_n)$, the
\emph{scatterplot} simply plots the points. It is the natural first step when we
suspect that $y$ depends on $x$. We will return to scatterplots in
Chapter~\ref{ch:leastsquares}, where we fit lines to them; see, for instance,
the timber and mortality data in that chapter.

\begin{example}[Correlation of geyser eruptions and waiting times]
For the geyser, each eruption's duration can be paired with the waiting time
until the next eruption. The empirical correlation coefficient of these two
quantities is about $0.90$—a strong positive association: long eruptions tend to
be followed by long waits.
\end{example}

\begin{example}[Height and weight, and a prediction]
For a classic dataset of women's heights and weights, fitting a line gives
$\text{height}\approx 0.287\,(\text{weight})+25.72$ inches, with an empirical
correlation of $0.995$—an exceptionally strong linear relationship. Using the
line, a woman weighing $105$ pounds is predicted to be about
$0.287(105)+25.72\approx 56$ inches tall.
\end{example}

\begin{example}[Linearizing exponential growth]
The U.S.\ population appears to grow roughly exponentially, $P(t)=ae^{kt}$.
Taking logarithms turns this into a \emph{linear} relationship,
$\ln P(t)=\ln a + kt$, so regressing $\ln P$ on $t$ recovers the parameters. With
$t$ measured in years from $1790$, this yields $\ln a\approx 1.688$ (so
$a\approx 5.41$ million) and $k\approx 0.0220$, i.e.\ $P(t)\approx 5.41\,e^{0.0220t}$.
Transforming to linearity before fitting is a recurring EDA tactic.
\end{example}

\section{Numerical Summaries}

Graphs reveal shape; numbers pin down location and spread.

\subsection{The center of a dataset}

The \emph{sample mean} is the ordinary average
$\xbar_n=\tfrac1n\sum_{i=1}^n x_i$, the natural analogue of expectation. The
\emph{sample median} is the middle value of the sorted data (the average of the
two middle values when $n$ is even). The mean is sensitive to \emph{outliers};
the median is robust.

\begin{example}[Outliers and a recording error]
Eleven hourly temperatures recorded one New Year's at Wick, Scotland, were
\[
  43,43,41,41,41,42,43,58,58,41,41.
\]
The two $58$'s sit far from the rest and inflate the mean to $44.7$; dropping
them lowers it to $41.8$. The median barely moves ($42$ versus $41$). In fact the
two anomalies arose from a switch of recording units; the corrected values are
near $41$, giving mean $41.5$ and median $42$. The median's robustness is exactly
its insensitivity to such contamination. (This is an argument for being
\emph{aware} of outliers, not for blindly deleting them.)
\end{example}

\subsection{The spread of a dataset}

The \emph{sample variance} and \emph{sample standard deviation} are
\[
  s_n^2=\frac{1}{n-1}\sum_{i=1}^n (x_i-\xbar_n)^2,\qquad s_n=\sqrt{s_n^2}.
\]
The divisor $n-1$ rather than $n$ will be justified in
Chapter~\ref{ch:estimation}, where we show it makes $s_n^2$ unbiased. Like the
mean, the standard deviation is sensitive to outliers: for the uncorrected Wick
data $s_n=6.62$, but only $0.97$ with the two anomalies removed.

A robust alternative is the \emph{median of absolute deviations},
\[
  \MAD(x_1,\dots,x_n)=\Med\bigl(|x_1-\Med_n|,\dots,|x_n-\Med_n|\bigr),
\]
which, like the median, is hardly affected by a few extreme values.

\subsection{Quantiles, quartiles, and the IQR}

The \emph{$p$th empirical quantile} $q_n(p)$ is, loosely, the value below which a
proportion $p$ of the data lie. Writing the sorted data (the \emph{order
statistics}) as $x_{(1)}\le\cdots\le x_{(n)}$, we take the $i$th order statistic
to be the $\tfrac{i}{n+1}$ quantile and interpolate linearly between them.

\begin{example}[Interpolating a percentile]
For the Wick data, the order statistics put $42$ at the $\tfrac{6}{12}=0.5$
quantile and $43$ at the $\tfrac{7}{12}$ quantile. The $55$th percentile lies
between these; linear interpolation between $(0.5,42)$ and $(0.583,43)$ gives
$q_n(0.55)\approx 42.6$.
\end{example}

The lower and upper \emph{quartiles} are $q_n(0.25)$ and $q_n(0.75)$, and the
\emph{interquartile range} is
\[
  \mathrm{IQR}=q_n(0.75)-q_n(0.25),
\]
the span of the middle half of the data—another robust measure of spread.
Together with the minimum, median, and maximum, the quartiles form the
\emph{five-number summary}.

\begin{example}[A five-number summary]
For $1000$ recorded earthquake magnitudes, the five-number summary is
\[
  \min = 4.00,\quad Q_1 = 4.30,\quad \text{median}=4.60,\quad Q_3=4.90,\quad \max=6.40,
\]
with sample mean $4.62$. The closeness of mean and median, and the symmetry of
the quartiles about the median, suggest a roughly symmetric bulk with a modest
right tail.
\end{example}

\subsection{The boxplot}

The \emph{boxplot} renders the five-number summary graphically
(Figure~\ref{fig:boxplot}). A box spans the lower to upper quartile—its height is
the IQR—with a line at the median. \emph{Whiskers} extend to the most extreme
observations within $1.5\times\mathrm{IQR}$ of the box; anything beyond is plotted
individually as a suspected outlier.

\begin{figure}[ht]
  \centering
  \includegraphics[width=\linewidth]{boxplot}
  \caption{Left: the anatomy of a boxplot. Right: a boxplot of the geyser data.
  Boxplots are excellent for spread and skew but, as here, cannot reveal that a
  dataset is bimodal—the off-center median is the only hint of structure.}
  \label{fig:boxplot}
\end{figure}

The boxplot is superb for comparing spreads and spotting skew and outliers, but
it has a blind spot: it cannot show that the geyser data are \emph{bimodal}. The
position of the median within the box hints at asymmetry, but only the histogram
and KDE reveal the two peaks. No single summary tells the whole story—which is
why exploratory analysis uses several in concert.
